% /chapters/2-descriptive-combinatorics/01-background.tex

\section{Background}\label{sec:descr_combi_background}

Throughout this entire chapter, $G$ denotes an analytic graph on a Polish space $V(G)=X$.
That is, $X$ is a separable completely metrizable topological space, and $G$ is a symmetric and irreflexive relation on $X$ which, as a subspace of $X\times X$, is the continuous image of a Polish space.

I employ the symbol $H\to_cG$ (respectively, $H\to_BG$) to abbreviate the fact that there is a continuous (respectively Borel) homomorphism from the analytic graph $H$ to the analytic graph $G$.

\subsection{The Borel chromatic number}

The \emph{Borel chromatic number} $\chi_B(G)$ of a Borel graph~$G$ was introduced by Kechris, Solecki, and Todorčević~\cite{kechris1999borel}, who proved the celebrated \emph{$G_0$ dichotomy}: there exists a single Borel graph~$\mathbb{G}_0$ such that
\begin{align}\label{eq:G0-dichotomy}
    \chi_B(G) > \aleph_0\quad
    \Longleftrightarrow\quad
    \mathbb{G}_0 \to_B G.
\end{align}
In other words, the graph~$\mathbb{G}_0$ forms a one-element \emph{basis} for the class of Borel graphs with uncountable Borel chromatic number.

This dichotomy breaks down when one moves from uncountable to merely infinite Borel chromatic number.
In \cite{Todorcevic_2021_ComplexProbBorel}, Todorčević and Vidnyánszky recently proved that the set of Borel graphs with infinite Borel chromatic number (equivalently, Borel chromatic number $\geq 4$ for closed subgraphs of the shift graph) is $\mathbf\Sigma^1_2$-complete.
In particular, no countable family of Borel graphs can serve as a basis for this class.
Moreover, their argument rules out the existence of a simple $G_0$-type dichotomy substituting $\aleph_0$ in \eqref{eq:G0-dichotomy} for any value in $3,4,5\dots$.
Interestingly, when substituting the value $2$, there again is a single Borel graph that makes a dichotomy like \eqref{eq:G0-dichotomy} hold.
This is known as the $L_0$~dichotomy and is the main result of \cite{L0paper}.
This thesis contains a more concise proof of said result in Section~\ref{sec:concise_L0}.

\begin{fact}\label{fact:chrom_transitive}
    If $H\to_cG$ or, more generally if $H\to_BG$, then $\chi_B(H)\leq\chi_B(G)$.
\end{fact}

Using a standard greedy algorithm argument, it is clear that any finite graph $G$ of maximum degree $\Delta$ satisfies $\chi(G)\leq\Delta+1$.
Interestingly, a similar fact holds for Borel graphs of uniformly bounded degree.

\begin{theorem}[Proposition~4.6 in \cite{kechris1999borel}]\label{prop:brooks}
    If $G$ is a Borel graph on a Polish space all of whose degrees are bounded by the natural number $k$, then $\chi_B(G)\leq k+1$.
\end{theorem}

A common situation is having an analytic set and the desire to apply to it a statement that holds for Borel sets instead.
One way to circumvent this problem is by using Borel separation to ``inflate'' the analytic set into a Borel set with the same property.
This technique is crucial, for instance, for the proof of the $G_0$ dichotomy found in \cite{Bernshteyn2024}.
Concretely, since the dichotomy is a statement about independent sets of vertices, the proof makes use of the following fact.

\begin{proposition}
    In the graph $G$, any analytic independent set of vertices is contained in a Borel independent set.
\end{proposition}

We present two proofs of this fact, original to this exposition, for several reasons.
First, because \cite{Bernshteyn2024} omits it~\footnote{The fact is indeed left as an exercise to the reader.}; it is an illustrative example of how different notions separation get used in an argument; it motivates the introduction of a more powerful technique that we need in Section~\ref{sec:concise_L0}.
Specifically, Claim~\ref{claim:separation_odd-walks} states an analog result for odd-walks, instead of merely independent sets.
The added complexity is that odd-walks can have different lengths, and so proving the aforementioned claim is not as straightforward without invoking Lemma~\ref{lemma:first-reflection}.

\begin{proof}
    Suppose that $I \subseteq V(G)$ is an analytic independent set.
    We view the edge relation of $G$ as a subset of $V(G)\times V(G)$, which is analytic by assumption.

    First, consider the set $$A_0 := \{x_0\in V(G):\exists x_1\in I\ (x_0x_1\in G)\},$$
    which consists of all vertices with at least one neighbour in $I$.
    To see that $A_0$ is analytic, write it as the projection onto the first coordinate of the set $$(V(G)\times I)\cap G \;=\; \{(x_0,x_1): x_1\in I \text{ and } x_0x_1\in G\}.$$
    This is the intersection of the analytic set $V(G)\times I$ with the analytic set $G$, hence analytic.
    Since $I$ is independent, $I\cap A_0=\emptyset$.
    As $I$ and $A_0$ are disjoint analytic sets, the Lusin Separation Theorem \cite[Theorem~14.7]{Kechris_1995_DescriptiveSetTheory} yields a Borel set $B_0\subseteq V(G)$ with $$I\subseteq B_0 \quad\text{and}\quad B_0\cap A_0=\emptyset.$$
    The second condition means that no vertex of $B_0$ has any neighbour in $I$.

    Next, let $$A_1 := \{x_1\in V(G):\exists x_0\in B_0\ (x_0x_1\in G)\}$$
    be the set of all vertices with at least one neighbour in $B_0$.
    Analogously to before, $A_1$ is the projection onto the second coordinate of
    $$(B_0\times V(G))\cap G \;=\; \{(x_0,x_1): x_0\in B_0 \text{ and } x_0x_1\in G\},$$
    which is the intersection of the Borel set $B_0\times V(G)$ with the analytic set $G$, and is therefore analytic.
    We claim that $I\cap A_1=\emptyset$.
    Indeed, if some $x_1\in I$ had a neighbour $x_0\in B_0$, then $x_0$ would be a vertex of $B_0$ with a neighbour $x_1\in I$, that is $x_0\in A_0$, contradicting $B_0\cap A_0=\emptyset$.
    So again we have two disjoint analytic sets, and separation yields a Borel set $B_1\subseteq V(G)$ with $$I\subseteq B_1 \quad\text{and}\quad B_1\cap A_1=\emptyset.$$
    The second condition means that no vertex of $B_1$ has any neighbour in $B_0$.

    Set $B:=B_0\cap B_1$, which is Borel and contains $I$.
    It remains to check that $B$ is independent.
    Suppose for contradiction that $x_0,x_1\in B$ with $x_0x_1\in G$.
    Then $x_0\in B_0$, so $x_1$ is a vertex with a neighbour in $B_0$, giving $x_1\in A_1$.
    But $x_1\in B_1$ and $B_1\cap A_1=\emptyset$, a contradiction.
\end{proof}

\begin{figure}[ht]
    \centering
    \input{figures/Borel-separation.tex}
    \caption{An illustration of Borel Separation in vertex sets.}
    \label{fig:Borel_separation}
\end{figure}

The following notion, taken from \cite[Theorem~35.10, p.~285]{Kechris_1995_DescriptiveSetTheory}, is a well-known fact about analytic sets.

\begin{definition}\label{def:pi11-on-sigma11}
    A collection of subsets $\Phi$ of a Polish space $X$ is said to be \emph{$\Pi^1_1$ on $\Sigma^1_1$} if for any Polish $Y$ and any $\Sigma^1_1$ set $A\subseteq Y\times X$, the set $A_\Phi:=\{y\in Y:A_y:=\{x\in X:(y,x)\in A\}\in\Phi\}$ is $\Pi^1_1$.
\end{definition}

\begin{lemma}[First Reflection Theorem]\label{lemma:first-reflection}
    If $\Phi$ is as in the preceding definition, then for any $\Sigma^1_1$ set $A\in\Phi$, there is a Borel set $B\subseteq X$ such that $A\subseteq B$ and $B\in\Phi$.
\end{lemma}

For example, the property of being an independent set of vertices is, as expected, $\Pi^1_1$ on $\Sigma^1_1$.
Indeed, let $\Phi$ denote the collection of all independent subsets of $V(G)$, let $Y$ be a Polish space, and let $A\subseteq Y\times V(G)$ be $\Sigma^1_1$.
We must show that $A_\Phi=\{y\in Y: A_y\in\Phi\}$ is $\Pi^1_1$.
Observe that its complement consists precisely of those $y\in Y$ for which $A_y$ contains an edge:
$$Y\setminus A_\Phi = \{y\in Y:\exists x_0,x_1\in V(G)\ (x_0x_1\in E(G)\wedge (y,x_0)\in A\wedge (y,x_1)\in A)\}.$$
The set $\{(y,x_0,x_1): (y,x_0)\in A\wedge (y,x_1)\in A\}$ is $\Sigma^1_1$, being the preimage of $A\times A$ under the continuous map $(y,x_0,x_1)\mapsto ((y,x_0),(y,x_1))$; intersecting with the Borel set $\{(y,x_0,x_1): x_0x_1\in E(G)\}$ keeps it $\Sigma^1_1$.
Projecting onto $Y$ therefore yields a $\Sigma^1_1$ set, so $A_\Phi$ is $\Pi^1_1$, as required.

Consequently, the Proposition above is an immediate corollary of Lemma~\ref{lemma:first-reflection}: given an analytic independent set $I\in\Phi$, the First Reflection Theorem directly produces a Borel set $B\supseteq I$ with $B\in\Phi$.
Note that this second proof, unlike the elementary one given above, generalises readily: one only needs to verify the $\Pi^1_1$ on $\Sigma^1_1$ property for the relevant family $\Phi$, which, for more complex combinatorial properties, is typically a non-trivial step.

\subsection{The Borel dichromatic number}

The \emph{dichromatic number} of a directed graph was introduced by Neumann-Lara~\cite{NeumannLara1982} as a natural directed analogue of the chromatic number.
A \emph{dicolouring} of a directed graph~$D$ is a map $c\colon V(D)\to k$ such that every colour class induces a subgraph without any directed cycles; the \emph{dichromatic number} $\vec\chi(D)$ is the least such~$k$.
Equivalently, $\vec\chi(D)$ is the minimum number of acyclic sets into which the vertex set can be partitioned.

This notion has attracted considerable attention in finite combinatorics.
Neumann-Lara~\cite{NeumannLara1982} established several foundational results, including analogues of classical lower bounds for the chromatic number.
Bokal, Fijavž, Juvan, Kayll, and Mohar~\cite{Bokal_2004_CircularChrom} later introduced the circular dichromatic number of a directed graph and proved that the decision problem of whether a directed graph $D$ satisfies $\vec\chi(D)\leq 2$ or not is NP-complete; this is in contrast with the usual chromatic number for (undirected) graphs, where one gets NP-completeness only after posing the question for three colours or more.

Very recently, Raghavan and Xiao~\cite{RaghavanXiao2024} established a dichotomy for the Borel dichromatic number that generalizes the $G_0$ dichotomy to directed graphs: they showed that a continuum-size family of directed graphs characterizes when a Borel directed graph has uncountable Borel dichromatic number.
This continuum-size basis consists of pairwise Borel-incomparable directed graphs, but it is unknown whether any smaller basis exists; in particular, the question of whether a \emph{countable} basis suffices is open.

One can interpret the existence of a small (i.e.\ singleton, or finite, or countable) basis as a sort of measure of complexity of the associated decision problem.
For instance, the existence of $\mathbb{G}_0$ says that deciding whether a Borel graph has uncountable Borel chromatic number is ``straightforward:'' one only has to verify one possible ill-behaved property, namely containing a homomorphic copy of $\mathbb{G}_0$, to give a positive answer; and in some sense, finding this homomorphic copy is the \emph{only} way of proving that a Borel graph has uncountable Borel chromatic number.
In contrast, deciding whether a Borel graph has infinite Borel chromatic number is a much more complex task, as no easily describable method exists for deciding this in terms of graph homomorphisms, and so different graphs probably require wildly different proving methods, if any.

The Borel dichromatic number can recover the Borel chromatic number in the following sense.

\begin{remark}\label{rmk:dichrom-recovers-chrom}
	Given an undirected graph $G$, let $\vec{G}$ denote the directed graph obtained by replacing each edge $\{u,v\}$ with two opposing arcs $(u,v)$ and $(v,u)$.
	Then $\vec\chi_B(\vec{G}) = \chi_B(G)$.
	Indeed, a subset of vertices is acyclic in $\vec{G}$ if and only if it is independent in~$G$: any directed cycle in $\vec{G}$ would in particular contain an arc $(u,v)$ with $\{u,v\}\in E(G)$, so a set containing both $u$ and $v$ is neither acyclic in $\vec{G}$ nor independent in~$G$; conversely, an independent set in~$G$ induces a directed graph in~$\vec{G}$ with no arcs at all, hence trivially acyclic.
	Moreover, the map $\mathrm{code}(G)\mapsto\mathrm{code}(\vec{G})$ is $\Delta^1_1$, since it is built from the edge set using products and unions (Lemma~\ref{lemma:toolkit}(1),(3)).
	It follows that the $\mathbf\Sigma^1_2$-completeness of Borel $k$-colouring of undirected graphs~\cite{Todorcevic_2021_ComplexProbBorel} implies the $\mathbf\Sigma^1_2$-completeness of Borel $k$-dicolouring of directed graphs, for every $k\geq 4$.
	The case $k=2$ is not covered by this reduction (since deciding 2-colourability of undirected graphs is trivial), and is precisely the content of Proposition~\ref{prop:main}.
\end{remark}

It is natural to ask whether the Todorčević--Vidnyánszky phenomenon also occurs in the directed setting at the remaining finite threshold, or if there is a directed analog for the $L_0$ dichotomy:

\begin{quote}
\emph{Is there a countable basis for the class of Borel directed graphs with Borel dichromatic number at least~$3$?}
\end{quote}

Section~\ref{sec:dichromatic} gives a negative answer, thereby completing the picture for all finite thresholds of both the Borel chromatic and dichromatic numbers.

The following table summarizes the basis landscape, contrasting the undirected and directed cases at different thresholds.

\medskip
\begin{center}
\begin{tabular}{lll}
	\hline
	\textbf{Class} & \textbf{Existence of small basis} & \textbf{Reference} \\
	\hline
	$\chi_B(G) > \aleph_0$ & Yes (size 1: $\mathbb{G}_0$) & \cite{kechris1999borel} \\
	$\chi_B(G) > 2$ & Yes (size 1: $\mathbb{L}_0$) & \cite{L0paper}, Section~\ref{sec:concise_L0} \\
	$\chi_B(G) > k$, $k\geq 3$ & No ($\mathbf\Pi^1_2$-complete) & \cite{Todorcevic_2021_ComplexProbBorel} \\
	\hline
	$\vec\chi_B(D) > \aleph_0$ & Yes (size $\mathfrak{c}$) & \cite{RaghavanXiao2024} \\
	$\vec\chi_B(D) > k$, $k\geq 2$ & No ($\mathbf\Pi^1_2$-complete) & Prop.~\ref{prop:main}, Rmk.~\ref{rmk:dichrom-recovers-chrom} \\
	\hline
\end{tabular}
\end{center}